
%%%%%%%%%%%%%%%%%%%%%%%%%%%%%%%%%%%%%%%%%%%%%%%%%%%%%%%%%%%%%%%%%%%%%
%%%%%%%%%%%%%%%%%%%%%%%%%%%%%%%%%%%%%%%%%%%%%%%%%%%%%%%%%%%%%%%%%%%%%
%%%%%%%%%%%%%%%%%%%%%%%%%%%%%%%%%%%%%%%%%%%%%%%%%%%%%%%%%%%%%%%%%%%%%
\section{Graded quivers with potential and dilation classes}\label{sec:graded-QP}

The main goals of this section are to develop the notion of dilation classes, cf.~\Cref{ssec:dilation_classes}, and establish that graded mutations of quivers with potential provide a bijection between dilation classes, cf.~\Cref{thm:G1-mutation}.

%%%%%%%%%%%%%%%%%%%%%%%%%%%%%%%%%%%%%%%%%%%%%%%%%%%%%%%%%%%%%%%%%%%%% 
\subsection{Completed path algebras and QPs}\label{ssec:prelim_pathalgebras_and_QPs}

Let us summarize the ingredients from \cite{DWZ} that we use in this section, referring to \cite{DWZ} for further details. Let $Q$ be a quiver with vertex set $Q_0$ and arrow set $Q_1$. Consider
\[
  R_Q:=\bigoplus_{i\in Q_0}\kk e_i,
  \qquad
  A_Q:=\bigoplus_{a\in Q_1}\kk a,
\]
where $A_Q$ is an $R_Q$-bimodule by letting the idempotents $e_i\in R_Q$ act via concatenation, where $e_i$ is considered as a trivial constant path at the vertex $i\in Q_0$.  The completed path algebra $\widehat{\kk Q}$ of the quiver $Q$ and its cyclic quotient $\widehat{\kk Q}_{\cyc}$ are
\[
  \widehat{\kk Q}:=\widehat T_{R_Q}(A_Q)
  =\prod_{n\ge 0}A_Q^{\otimes_{R_Q}n},\qquad
  \widehat{\kk Q}_{\cyc}
  :=\widehat{\kk Q}/\overline{[\widehat{\kk Q},\widehat{\kk Q}]}.
\]
They respectively contain the ideals
\[
  \m_Q:=\prod_{n\ge1}A_Q^{\otimes_{R_Q}n}\sse \widehat{\kk Q},\qquad \m^2_{\cyc}\sse \widehat{\kk Q}_{\cyc}.
\]

\noindent The algebra morphisms that we will consider are all continuous for the \(\m_Q\)-adic topology. The following is \cite[Definition 4.1]{DWZ}:

\begin{definition}[Quiver with potential]
A potential $W$ for a quiver $Q$ is an element $W\in\m^2_{\cyc}$, and a quiver with potential (QP) is a pair $(Q,W)$.
\hfill$\Box$\end{definition}

\noindent The foundational aspects of the study of quivers with potential are developed in \cite{DWZ}. Quivers with potential will be considered up to right-equivalence, cf.~\cite[Definition 4.2]{DWZ}:

\begin{definition}[Right-equivalence]\label{def:right_equivalent}
Let $(Q,W)$ and $(Q',W')$ be two QPs on the same vertex set. A right-equivalence between $(Q,W)$ and $(Q',W')$ is a $\kk$-algebra isomorphism $\varphi:\widehat{\kk Q}\lr\widehat{\kk Q'}$ such that $\varphi|_{R_Q}=\mbox{id}$ and $\varphi(W)$ is cyclically equivalent to $W'$.\hfill$\Box$
\end{definition}

\noindent In Definition \ref{def:right_equivalent}, two potentials $W,W'$ are said to be cyclically equivalent if their difference $W-W'$ lies in the closure of the span of all elements of the form $a_1\cdots a_d-a_2\cdots a_da_1$, where $a_1,\ldots,a_d$ is a cyclic path, cf.~\cite[Definition 3.2]{DWZ}. We often denote a right-equivalence $\varphi$ from $(Q,W)$ to $(Q',W')$ by $\varphi:(Q,W)\lr(Q',W')$.\\

In \cite{DWZ} the notions of {\it reduced} and {\it trivial} QPs are also introduced, as follows. A QP $(Q,W)$ is said to be reduced if the degree-2 homogeneous part $W^{(2)}$ vanishes, i.e.~$W^{(2)}=0$. That is, $(Q,W)$ is reduced if $W$ contains no quadratic monomial terms, i.e.~no terms of the form $ab$, $a,b\in Q_1$. A QP $(Q,W)$ is said to be trivial if $W$ is entirely quadratic, i.e.~ $W\in \widehat{\kk Q}^{(2)}$ belongs to the degree-2 homogeneous part of the path algebra $\widehat{\kk Q}$, and the Jacobian algebra of $(Q,W)$ is isomorphic to the semisimple vertex algebra $R_Q$.

\begin{example}\label{ex:trivialQP} Let $Q$ be a quiver with two vertices $i,j\in Q_0$, $n\in\N$, and choose $2n$ arrows $x_\nu:i\to j$, $y_\nu:j\to i$, for $\nu\in[1,n]$. If we choose the potential $W=x_1y_1+\ldots+x_ny_n$, then $(Q,W)$ is a trivial QP.\hfill$\Box$
\end{example}

\noindent The Splitting Theorem \cite[Theorem~4.6]{DWZ} decomposes a QP, uniquely up to right-equivalence, as a direct sum of a reduced QP and a trivial quadratic QP, the latter as in \Cref{ex:trivialQP}.

\begin{definition}[Stable QP class]\label{def:stable-class}
The stable right-equivalence class $[Q,W]$ of a QP \((Q,W)\) is the set of QPs which can be obtained from $(Q,W)$ via a sequence intertwining right-equivalences, and adjoining or deleting trivial quadratic QP summands.
\hfill$\Box$\end{definition}

\noindent Note that QPs representing the same stable right-equivalence class have the same vertex set, up to relabeling. The notion of mutation of a quiver with potential consists of a pre-mutation followed by reduction to the reduced QP part, cf.~\cite[Section 5]{DWZ}. It is an involutive operation up to right-equivalence by \cite[Theorem~5.7]{DWZ}. We will also use the notion of a non-degenerate QP $(Q,W)$: a reduced QP is said to be non-degenerate if every finite sequence of QP mutations ends at a QP with a $2$-acyclic reduced underlying quiver, cf.~\cite[Section 7]{DWZ}.

%%%%%%%%%%%%%%%%%%%%%%%%%%%%%%%%%%%%%%%%%%%%%%%%%%%%%%%%%%%%%%%%%%%%%
\subsection{$\Z$-graded QPs and vertex gauge}

We use the notion of $\Z$-graded QPs and their mutations, as introduced in \cite[Section 6]{AO}; see also \cite{Keller2026_GradedPreprint} and \cite{HerschendIyama2010}. We must add a new notion beyond those in \cite{AO}, that of a vertex gauge. This subsection presents the necessary ingredients on $\Z$-graded QPs and vertex gauge. The following definition appears in \cite[Section 6.3]{AO}:

\begin{definition}[Degree-one $\Z$-grading]\label{def:degree-one-grading}
Let $(Q,W)$ be a QP. By definition, a $\Z$-grading of $(Q,W)$ is a function
\[
 d:Q_1\longrightarrow\Z.
\]
By extending additively to paths, $W$ is said to be of degree one with respect to $d$ if every homogeneous component of \(W\) has total internal degree
one. Equivalently, \(W\) is homogeneous of degree one in the completed
cyclic quotient.
\hfill$\Box$\end{definition}
%\noindent Note that an arrow $\Z$-grading on a quiver $Q$ naturally enhances the path algebra $\widehat{\kk Q}$ to a $\Z$-graded algebra.
\noindent In this manuscript we only consider $\Z$-graded QPs of degree one: to ease notation, a graded QP $(Q,W,d)$ is henceforth understood to be a QP $(Q,W)$ endowed with a degree-one $\Z$-grading as in \Cref{def:degree-one-grading}. Note that a $\Z$-grading $d:Q_1\lr\Z$ for which $W$ is homogeneous of degree one and $d$ only takes value in $\{0,1\}\sse\Z$ relates to the notion of a cut, cf.~\cite[Definition 3.1]{HerschendIyama2010}.\\

Note that \Cref{def:degree-one-grading} grades the arrows of $Q$, rather than its vertices. Grading vertices yields an equivalence relation on arrow $\Z$-gradings, as follows:

\begin{definition}[Vertex gauge]\label{def:vertex-gauge}
Let $(Q,W,d)$ be a graded QP. Given a vertex $\Z$-grading \(h:Q_0\to\Z\) of $Q$, consider the arrow $\Z$-grading
\begin{equation}\label{eq:gauge}
 d^h(a):=d(a)+h(t(a))-h(s(a)).
\end{equation}
By definition, the two graded QPs $(Q,W,d)$ and $(Q,W,d^h)$ are said to be vertex-gauge equivalent to each other, i.e.~two graded QPs are vertex-gauge equivalent if one is obtained from the
other by changing the $\Z$-grading via \eqref{eq:gauge}.
\hfill$\Box$\end{definition}

\noindent In \Cref{def:vertex-gauge}, $s(a),t(a)\in Q_0$ denote the source and target vertices of an arrow $a\in Q_1$. The notion of a vertex gauge is well-defined since vertex-gauging preserves the degree-one property of $W$. Indeed, suppose that $W$ is of degree one with respect to $d$. Then for a cyclic monomial \(a_m\cdots a_1\) in $W$, the difference between the $d$ and $d^h$ degrees of the monomial $a_m\cdots a_1$ is
\[
 \sum_{r=1}^m\bigl(h(t(a_r))-h(s(a_r))\bigr)=0,
\]
since the sum telescopes around the cycle.  Hence every cyclic monomial has the same
degree for \(d\) and \(d^h\) and \Cref{def:vertex-gauge} is well-defined.\\

%%%%%%%%%%%%%%%%%%%%%%%%%%%%%%%%%%%%%%%%%%%%%%%%%%%%%%%%%%%%%%%%%%%%% 
\subsection{Graded right-equivalence and homogeneous stabilization}

In the same manner that QPs are considered up to right-equivalence and stabilizations (i.e.~adjoining or deleting a trivial quadratic QP summand), and there is a Splitting Theorem, these notions also exist for graded QPs, as explained in \cite[Section 6.3]{AO}. For completeness, we recall them, see also \cite{Keller2026_GradedPreprint}. First, graded right-equivalence is \cite[Definition~6.5]{AO}:

\begin{definition}[Graded right-equivalence]\label{def:graded-right-equivalence}
Let $(Q,W,d),(Q',W',d')$ be two graded QPs. A graded right-equivalence $\varphi:\widehat{\kk Q}\lr\widehat{\kk Q'}$ is a right-equivalence from $(Q,W)$ to $(Q',W')$ such that $(\widehat{\kk Q},d)\lr(\widehat{\kk Q'},d')$ is an isomorphism of graded algebras.
\hfill$\Box$\end{definition}

\begin{definition}[Homogeneous stabilization]\label{def:hom-trivial}
By definition, a graded QP $(Q,W,d)$ is said to be trivial if the potential $W$ is entirely quadratic, homogeneous of $d$-degree one, and its Jacobian algebra is isomorphic to the vertex algebra $R_Q$. By definition, two graded QPs are said to be related by a homogeneous stabilization or reduction if one is obtained from the other by adjoining or deleting a trivial graded QP.
\hfill$\Box$\end{definition}

The graded version of the Splitting Theorem \cite[Theorem~4.6]{DWZ} is \cite[Theorem~6.6]{AO}: it states that a graded QP splits into reduced and
trivial graded parts, each unique up to graded right-equivalence. See also \cite{Keller2026_GradedPreprint} for more details.

%%%%%%%%%%%%%%%%%%%%%%%%%%%%%%%%%%%%%%%%%%%%%%%%%%%%%%%%%%%%%%%%%%%%% 
\subsection{Dilation classes for QPs}\label{ssec:dilation_classes}

Let us now introduce a set that encodes the possible gradings of a QP $(Q,W)$ and incorporates the necessary equivalence relations to be able to work invariantly:

\begin{definition}[Dilations]\label{def:G1}
Let $[Q,W]$ be a stable right-equivalence class of a QP $(Q,W)$. By definition, the set of dilations $\Gclass[Q,W]$ associated to $[Q,W]$ is the set of all equivalence classes of graded QPs \((Q',W',d')\) such that $[Q',W']=[Q,W]$, where the equivalence relation is generated by the following:
\begin{enumerate}[label=(\alph*)]
\item vertex gauge, as in \Cref{def:vertex-gauge},
\item graded right-equivalence, as in \Cref{def:graded-right-equivalence},
\item homogeneous trivial stabilization and reduction, as in \Cref{def:hom-trivial}.
\end{enumerate}
By definition, an element \(\delta\in\Gclass[Q,W]\) is said to be a dilation or a dilation class.
\hfill$\Box$\end{definition}

\noindent By \Cref{def:G1}, the quiver $Q'$ of any representative of a dilation class $\delta=[(Q',W',d')]\in\Gclass[Q,W]$ has a vertex set $Q'_0$ identical to the vertex set $Q_0$ of Q, as neither right-equivalence nor stabilization of QPs, and none of the three equivalence relations (a), (b), (c) in \Cref{def:G1}, affect the vertex set of $Q$. Similarly, the arrow set $Q'_1$ of any such quiver differs from the arrow set $Q_1$ of $Q$ solely in adding or subtracting canceling pairs of arrows.

\begin{remark} If one considers $Q$ up to relabeling of the vertices, one can incorporate simultaneous vertex relabeling as a fourth equivalence relation to be imposed in \Cref{def:G1}. The main results of the article are independent of vertex relabeling.\hfill$\Box$
\end{remark}

\noindent Note that a right-equivalence of (ungraded) QPs from $(Q,W)$ to $(Q',W')$ might not be compatible with some given gradings $(Q,W,d)$ to $(Q',W',d')$ on those QPs. That is, the dilations associated to $(Q,W,d)$ and $(Q',W',d')$, where $d,d'$ are chosen independently of each other, are typically different in $\Gclass[Q,W]$ even if $(Q,W)$ and $(Q',W')$ are right-equivalent (or even identical). That said, given a right-equivalence $\varphi$ from $(Q,W)$ to $(Q',W')$ and a grading $(Q,W,d)$, it is possible to transport $d$ via $\varphi$ to a grading of $(Q',W')$, as follows.

\begin{lemma}[Homogeneous rectification]\label{lem:rectification}
Let \((Q,W,d)\) be a graded QP, $(Q',W')$ an ungraded QP, and $\varphi:(Q,W)\lr(Q',W')$ an $($ungraded$)$ right-equivalence. Then the following three objects exist:
\begin{enumerate}
\item a QP $(Q',V)$ right-equivalent to $(Q',W')$,
\item a graded QP $(Q',V,d')$,
\item a graded right-equivalence
\[
 \psi:(Q,W,d)\longrightarrow(Q',V,d').
\]
\end{enumerate}
In addition, the dilation \([(Q',V,d')]\in\Gclass[Q',W']\) is
independent of all choices.
\end{lemma}

\begin{proof}
First, note that the arrow grading \(d\) gives a complete internal \(\Z\)-grading on
\(\widehat{\kk Q}\): each path is homogeneous, with degree the sum of its arrow
degrees, and an arbitrary formal series is resolved coefficient-wise into its
homogeneous parts. We transport this grading through the (ungraded) right-equivalence \(\varphi\) to
\(\widehat{\kk Q'}\) by declaring
\[
 \widehat{\kk Q'}_r:=\varphi(\widehat{\kk Q}_r),
\]
where the subindex $r$ indicates the piece with internal degree $r$. Because \(\varphi\) fixes the vertex algebra and preserves the arrow ideal $\mathfrak{m}_Q$,
\(R_{Q'}\) is concentrated in degree zero, \(\m_{Q'}\) is homogeneous, and
\(\m_{Q'}/\m_{Q'}^2\) is a finite-dimensional graded
\(R_{Q'}\)-bimodule. Now choose a homogeneous \(R_{Q'}\)-bimodule section
\[
 s:\m_{Q'}/\m_{Q'}^2\longrightarrow\m_{Q'}.
\]
Such a section exists: e.g.~choose a linear section separately in every
internal degree and every \(e_jA_{Q'}e_i\).  Choose a homogeneous basis
of \(\m_{Q'}/\m_{Q'}^2\) and identify it with the displayed arrow space
\(A_{Q'}\).  Let \(d'\) be the resulting arrow grading.  By the universal
property of the completed tensor algebra, the assignment sending each
displayed arrow to its lift under the section \(s\) extends to a continuous endomorphism
\[
 \theta:\widehat{\kk Q'}\longrightarrow \widehat{\kk Q'},
\]
which fixes the idempotent elements. Its linear part on \(\m_{Q'}/\m_{Q'}^2\) is invertible, so the formal inverse
lemma applied to these complete path algebras shows that \(\theta\) is an automorphism. We then define
\[
 V:=\theta^{-1}(W'),
 \qquad
 \psi:=\theta^{-1}\circ\varphi.
\]
By construction, the map \(\psi\) is graded from \(d\) to \(d'\), and
\(\psi(W)=V\) in the cyclic quotient.  Hence \(V\) is homogeneous of degree
one and \(\psi\) is a graded right-equivalence. This establishes the existence of the objects $(1),(2),(3)$.\\

For independence of choices, let \(\theta_1,\theta_2\) be two rectifications, as above, producing
\((V_1,d_1')\) and \((V_2,d_2')\).  Since both identify their corresponding graded
algebras with the same transported graded algebra \(\widehat{\kk Q'}\) we conclude that
\[
 \theta_2^{-1}\theta_1:(Q',V_1,d_1')\longrightarrow(Q',V_2,d_2')
\]
is a graded right-equivalence, sending $V_1$ to $V_2$ in the corresponding cyclic quotients.
\end{proof}
\color{black}
Since homogeneous stabilizations provide a bijection between dilation classes, it then follows from the homogeneous rectification provided by \Cref{lem:rectification} that we have a canonical bijection between the sets of dilations associated to two QPs that are right-equivalent:

\begin{corollary}\label{cor:G1-presentation}
If \((Q,W)\) and \((Q',W')\) are stably right-equivalent, then homogeneous
rectification induces a canonical bijection
\[
 \Gclass[Q,W]
 \xrightarrow{\sim}
 \Gclass[Q',W']
\]
between their sets of dilations.\hfill$\Box$
\end{corollary}

%%%%%%%%%%%%%%%%%%%%%%%%%%%%%%%%%%%%%%%%%%%%%%%%%%%%%%%%%%%%%%%%%%%%%
%%%%%%%%%%%%%%%%%%%%%%%%%%%%%%%%%%%%%%%%%%%%%%%%%%%%%%%%%%%%%%%%%%%%%
%%%%%%%%%%%%%%%%%%%%%%%%%%%%%%%%%%%%%%%%%%%%%%%%%%%%%%%%%%%%%%%%%%%%%
\subsection{Graded QP mutation and vertex gauge}\label{sec:graded-mutation}

The QP mutation of \cite[Section 5]{DWZ} is generalized to a graded QP mutation in \cite[Section 6.3]{AO}, and see also \cite{Keller2026_GradedPreprint}. After briefly recalling such notion, we show in \Cref{thm:G1-mutation} that graded QP mutation induces a canonical bijection between the dilations of a QP $(Q,W)$ and its QP-mutation $\mu_k(Q,W)$ at a vertex $k$. Let \((Q,W,d)\) be a graded QP and \(k\in Q_0\) be a vertex incident
to no loop or oriented two-cycle. Let us also choose a cyclic representative of $W$ so that no monomial begins or ends at \(k\).\\

The graded QP mutation of \cite[Section 6.3]{AO} consists of two steps, as does QP-mutation: a pre-mutation and a reduction, the latter using the corresponding Splitting Theorem. The pre-mutation step is as follows. Given two arrows
\[
 a:i\longrightarrow k,
 \qquad
 b:k\longrightarrow j,\qquad i,j,k\in Q_0,\quad a,b\in Q_1,
\]
the QP pre-mutation step creates a composite \([ba]:i\to j\), reverses \(a,b\), and
adds the corresponding cubic term to the potential $W$, resulting in a QP $(\tilde{\mu}_k^L Q,\tilde{\mu}_k^L W)$. By definition, the {\it left graded} pre-mutation assigns gradings
\begin{align}
 \widetilde d^L([ba])&=d(b)+d(a),\label{eq:mut-composite}\\
 \widetilde d^L(a^*)&=1-d(a),\label{eq:mut-incoming}\\
 \widetilde d^L(b^*)&=-d(b),\label{eq:mut-outgoing}
\end{align}
while unchanged arrows retain their degree. The pre-mutated potential $\tilde{\mu}_k^LW$ is
homogeneous of degree one with respect to this pre-mutated grading $\tilde{\mu}_k^L(d):=\widetilde{d}^L$. This thus defines a new graded QP
$$\tilde{\mu}_k^L(Q,W,d):=(\tilde{\mu}_k^L Q,\tilde{\mu}_k^L W,\tilde{\mu}_k^L d).$$
Then \cite[Definition 6.8]{AO}, which uses \cite[Thm.~6.6 \& Prop.~6.7]{AO}, reads as follows:

\begin{definition}[Left graded QP-mutation]
Let $(Q,W,d)$ be a graded QP and $k\in Q_0$ as above. By definition, its left graded QP-mutation
$$\mu_k^L(Q,W,d)=(\mu_k^LQ,\mu_k^LW,\mu_k^Ld)$$
is the graded QP given by the reduced graded part of the pre-mutated graded QP $\tilde{\mu}_k^L(Q,W,d)$.\hfill$\Box$
\end{definition}

\begin{remark}
There is an analogous notion of right graded QP-mutation $\mu_k^R(Q,W,d)$: the only change is that the pre-mutated degree $d_R$ assigns $d_R(a^*)=-d(a)$ and $d_R(b^*)=1-d(b)$.\hfill$\Box$
\end{remark}

Before establishing that both left and right graded QP-mutations induce bijections between the sets of dilations, we must show that the vertex gauging in \Cref{def:vertex-gauge} commutes with graded QP-mutations. The statement reads as follows:

\begin{lemma}[Graded QP-mutation commutes with vertex gauge]\label{prop:gauge-mutation}
Let $(Q,W,d)$ be a graded QP and \(h:Q_0\to\Z\) a vertex grading. Then:

\begin{enumerate}
    \item We have an equality of (arrow) gradings
    \begin{equation}\label{eq:gauge-premutation}
 \widetilde\mu_k^L(d^h)
 =\bigl(\widetilde\mu_k^L(d)\bigr)^h,
\end{equation}
where $h$ is understood as the same function \(h\) on the mutated vertex set, and the right hand side of \Cref{eq:gauge-premutation} denotes applying the corresponding vertex gauge to $\widetilde\mu_k^L(d)$.\\

 \item The left-graded mutated QP satisfies
 \begin{equation}\label{eq:gauge-reduced-mutation}
 \mu_k^L(Q,W,d^h)
 \simeq
 \mu_k^L(Q,W,d)^h,
\end{equation}
where $\simeq$ in \Cref{eq:gauge-reduced-mutation} denotes graded right-equivalence, and its right hand side denotes applying the corresponding vertex gauge to $\mu_k^L(d)$.\\

\item Both statements above hold for right graded QP-mutations.

\end{enumerate}
\end{lemma}

\begin{proof}
For Part (1), consider the composite arrow \([ba]:i\to j\). Then
\begin{align*}
 \widetilde d^{\,h}([ba])
 &=d^h(b)+d^h(a)\\
 &=d(b)+h(j)-h(k)+d(a)+h(k)-h(i)\\
 &=\widetilde d([ba])+h(j)-h(i).
\end{align*}
Similarly, for \(a^*:k\to i\) we have
\[
 \widetilde d^{\,h}(a^*)
 =1-d^h(a)
 =1-d(a)+h(i)-h(k),
\]
and for \(b^*:j\to k\) we obtain
\[
 \widetilde d^{\,h}(b^*)
 =-d^h(b)
 =-d(b)+h(k)-h(j).
\]
These are exactly the vertex-gauge corrections for the mutated arrows, thus proving
\Cref{eq:gauge-premutation}. For Part (2), a graded splitting right-equivalence remains
graded after applying the same vertex-gauge, because every path from \(i\) to \(j\)
has its degree shifted by the common amount \(h(j)-h(i)\). Hence graded
reduction and \Cref{eq:gauge-premutation} together imply \Cref{eq:gauge-reduced-mutation}. Part (3) is proven analogously.
\end{proof}

\begin{theorem}[Graded QP-mutation gives bijection on dilation classes]\label{thm:G1-mutation}
Let $[Q,W]$ be a stable right-equivalence class of a reduced
non-degenerate QP $(Q,W)$ and let \(k\) be a vertex. Then left graded QP-mutation induces a
canonical bijection
\[
 \mu_k^L:\Gclass[Q,W]
 \xrightarrow{\sim}
 \Gclass[\mu_k(Q,W)].
\]
Furthermore, right graded QP-mutation provides its inverse.
\end{theorem}

\begin{proof}
We must verify the three equivalence relations from \Cref{def:G1}. For vertex gauge, \Cref{prop:gauge-mutation} gives compatibility with vertex gauge. For graded right-equivalence: graded pre-mutation is functorial under graded
right-equivalence, and the Graded Splitting Theorem \cite[Theorem 6.6]{AO} makes the reduced output
well-defined up to graded right-equivalence. Similarly, homogeneous stabilization is absorbed by the Graded Splitting Theorem. Therefore both left and
right graded QP-mutations descend to maps on dilation classes. Since \cite[Lemma 6.9]{AO} shows that
\[
 \mu_k^R\mu_k^L(Q,W,d)
 \simeq(Q,W,d),
 \qquad
 \mu_k^L\mu_k^R(Q,W,d)
 \simeq(Q,W,d),
\]
where $\simeq$ denotes a graded right-equivalence, the induced maps on the set of dilations must therefore be inverse bijections.
\end{proof}